%% Abstract text only; layout is set by \hxmaketitle in main.tex
Multirotor UAVs lose their GNSS position indoors, underground, under tree canopy and
when signals are jammed, and without an absolute fix inertial navigation drifts
quadratically. We present a low-cost, open-source hardware and software system that
gives a Pixhawk-class multirotor running ArduPilot an independent position and obstacle
reference. A Raspberry Pi on the vehicle streams autopilot telemetry (including optical
flow and rangefinder), 360\textdegree{} YDLIDAR~G2 scans and camera video over Wi-Fi
to a ground laptop running ROS~2 Humble. There, odometry is reconstructed, Nav2 plans
collision-free velocity commands, and a hybrid perception node refines YOLOv8 masks
with Segment Anything~2.1. Commands return through a gated MAVLink path that acts only
in armed \code{GUIDED} mode, is clamped to \SI{0.5}{\metre\per\second} and stops after
\SI{0.5}{\second} without input. Our link optimisations cut command forwarding from
\SI{54}{\milli\second} to \SI{1.4}{\milli\second} and telemetry-to-odometry delay from
\SI{20}{\milli\second} to \SI{2}{\milli\second} (single host). We validated the system with a
bench log of sensor noise and the zero-velocity detector, a dead-reckoning
experiment (\SI{86.5}{\metre} error after \SI{25}{\second} uncompensated), a GPU
perception benchmark (\SI{5.6}{\milli\second} YOLO-only, \SI{31.6}{\milli\second} with
SAM), 17/17 passing unit tests, and 26 simulated flights (\SI{1.7}{h}) against Gazebo ground truth: speed
limits held, height hold was within 3--\SI{14}{cm} RMS, and three explorations
completed, while four office flights fell after propeller contact. On the real vehicle,
Nav2 flew an autonomous square path around trees outdoors at night, and a
\SI{7}{\minute} GNSS-free hover at \SI{2.5}{\metre} over grass drifted \SI{1.6}{\metre}, bounded but over 30 times the
\SI{5}{\centi\metre} target; trajectory accuracy against ground truth is not yet
measured. We also describe a fully onboard stereo-vision design
(\INR{45{,}679} bill of materials) as the next step.
